\documentclass[../main.tex]{subfiles}
\begin{document}
%==========================================TIÊU ĐỀ TẬP========================================
\begin{table}[h]
    \begin{adjustwidth}{-1cm}{}
        \begin{tabular}{>{\centering\arraybackslash}p{6cm}|>{\raggedright\arraybackslash}p{12.5cm}}
            \multirow{5}{6cm}
            {\\[-40pt]
                \flaregothic\fontsize{20pt}{20pt}\selectfont \centering
                \color{eptype}HISTOIRE\color{black}\\
                \freshbot\fontsize{72pt}{72pt}\selectfont
                \color{epnum}23\\[6pt]
                \cafeteria\fontsize{20pt}{20pt}\selectfont\color{epnum}(D-23)\color{black}
                \\[18pt]
                \color{black}
            }
             &
            {\fotskip\fontsize{18pt}{18pt}\selectfont \ruby{再}{さい}\ruby{会}{かい}への\ruby{電}{でん}\ruby{話}{わ}}
            \\[6pt]
             & {\cafeteria\fontsize{18pt}{18pt}\selectfont A Call About Reunion}                        \\[4pt]
             & {\vnmsans\fontsize{18pt}{18pt}\selectfont Cuộc gọi về chiến dịch Réunion}                \\[8pt]
             & {\caviar{\fontsize{14pt}{14pt}\selectfont Nhân vật: \emp{kuon2} \emp{achan} \emp{game}}} \\
             & {\abv Truyện phụ:} \emp{delu}
            \\[8pt]
        \end{tabular}
    \end{adjustwidth}
\end{table}
%=========================================TRUYỆN CHÍNH========================================
\color{black}

\txtbx[2]{
    {\color{vol} \uline{Tóm tắt:}} Sau đêm chiếu phim, A-chan gọi riêng cho Kuon và tiết lộ kế hoạch mang tên ``chiến dịch Réunion'' nhằm tạo điều kiện để các thành viên còn hoạt động gặp lại những người đã tốt nghiệp, trong một không gian riêng tư, có sự đồng thuận của tất cả và không gắn với lời hứa quay lại công ty. Cuộc gọi khiến Kuon vừa hy vọng vừa dè dặt, nhất là khi người đầu tiên được nhắc đến là Mano Aloe, nay là Delutaya. Cậu đồng ý làm đầu mối hỗ trợ, với điều kiện mọi bước đi phải tôn trọng lựa chọn và sự an toàn của cô. Một tin nhắn từ Nene mở ra bước đầu tiên cho cuộc đoàn tụ holoFive.
}

\pov{Hikawa Kuon}

Ba ngày trôi qua kể từ buổi chiếu phim chung của mọi người ở ký túc xá. Tôi vừa đóng cửa văn phòng sau một buổi họp muộn, chiếc điện thoại trong túi áo đột nhiên rung lên một nhịp đều đặn, phá vỡ sự yên tĩnh của hành lang.

Tôi rút máy ra, ánh sáng màn hình chiếu thẳng vào mặt. Tên người gọi hiện lên rõ ràng: Yuujin A-san. Một cuộc gọi riêng từ người đồng nghiệp vào giờ này, không phải là điều thường xảy ra. Tôi dừng bước ngay tại khoảng hành lang trống trải cạnh phòng họp lớn, nơi hiếm khi có người qua lại sau khi giờ làm việc chính thức kết thúc, rồi mới nhấn nút nghe.

``Hikawa Kuon xin nghe đây ạ.''

Đầu dây bên kia im lặng trong vài giây, chỉ có tiếng hơi thở nhẹ thoảng qua loa. Tôi có thể cảm nhận được sự chần chừ nào đó từ phía chị A, trước khi giọng nói quen thuộc vang lên, nhỏ hơn hẳn so với những lần chúng ta trao đổi công việc thông thường, như thể chị cũng đang ở một nơi nào đó vắng vẻ và không muốn ai nghe thấy cuộc hội thoại này: ``Kuon-san, em đang ở một mình chứ?''

Tôi quay người nhìn dọc theo toàn bộ hành lang. Các cánh cửa văn phòng khác đều đã đóng kín, đèn chỉ còn sáng ở mấy cái đèn trần cách xa nhau, không có bóng người nào di chuyển hay đứng ở gần. Không có ai có thể tình cờ nghe được nội dung cuộc gọi nếu tôi giữ giọng ở mức vừa phải.

``Vâng. Em đang ở một mình. Có chuyện gì vậy chị?''

``Chị muốn hỏi thăm tình hình ở ký túc xá. Mấy ngày nay mọi người vẫn ổn chứ?''

Tôi dựa lưng vào bức tường sơn trắng, mắt liếc qua ô cửa kính lớn ở cuối hành lang, nơi có thể nhìn thấy toàn bộ bãi đỗ xe dưới tầng trệt. Ánh sáng vàng từ đèn đường hắt vào, tạo ra những bóng dài vắt ngang trên sàn đá.

``Vâng. Nhà vẫn ổn ạ. Mọi người đều bình thường.''

``Vậy có chuyện gì bất thường nào xảy ra không? Hay có ai có biểu hiện gì khác thường không?''

Tôi không trả lời ngay lập tức. Câu hỏi của A-san làm những hình ảnh từ đêm chiếu phim ùa về trong đầu, như những thước phim cũ được tua lại. Tôi nhớ lại Mikeneko đứng tự bao lâu ở ngưỡng cửa phòng khách, bóng người cô dường như muốn hòa vào bóng tối, không bước vào mà cũng không rời đi. Rồi là Koneko, đã ngủ gục trên vai Naomi từ giữa bộ phim, miệng còn cắn một miếng bỏng ngô chưa kịp nhai.

Và cuối cùng, hình ảnh Delutaya ngồi ở góc xa nhất của ghế sofa, không nói một lời từ đầu đến cuối, ánh mắt chỉ dán vào màn hình chiếu lớn, ngay cả khi những cảnh cũ về \hll\ thế hệ thứ năm xuất hiện. Cô ấy cũng không rủ ai nói chuyện, chỉ lẳng lặng rời phòng khi phim kết thúc.

``Không có gì cần báo cáo ngay lúc này,'' tôi cuối cùng cũng nói ra câu đó.

Đó không hẳn là một lời nói dối hoàn toàn. Không có ai bị đau ốm, không có xung đột nào nảy ra, mọi sinh hoạt trong nhà vẫn diễn ra như bình thường. Nhưng nó cũng không phải là toàn bộ sự thật. Những thay đổi nhỏ trong tâm trạng của từng người, những khoảng lặng khó tả trong không khí nhà, những ánh mắt né tránh khi một cái tên cũ được nhắc đến\dots\ tất cả những điều đó không thể nào viết vào một bản báo cáo, hay nói ra qua điện thoại chỉ trong vài câu.

A-san dường như ngay lập tức nhận ra sự ngập ngừng trong giọng nói của tôi, sự trì hoãn trước khi tôi trả lời. Chị thở nhẹ ở đầu dây, như thể đã đoán trước được phản ứng của tôi.

``Chị hiểu. Chị không định hỏi về những chuyện riêng tư của từng người ở trong nhà. Chị biết em luôn giữ gìn không gian riêng cho mọi người.''

Tôi im lặng, chờ chị nói tiếp. Tôi biết nếu chỉ để hỏi thăm tình hình thông thường, chị sẽ không gọi điện riêng vào giờ này, với giọng nói cẩn trọng đến vậy. Chắc chắn có một việc quan trọng hơn đang ở phía sau.

``Chị gọi vì có một việc muốn trao đổi thẳng thắn với em. Một kế hoạch mới của công ty, chị muốn hỏi ý kiến em trước khi tiến hành bất cứ bước nào.''

``\dots Kế hoạch này có liên quan đến ký túc xá ạ?'' tôi lập tức hỏi ngay. Nếu nó liên quan đến những người đang sống trong nhà, tôi cần phải biết ngay từ đầu.

``Có thể nói là vậy. Nó liên quan đến những người đã từng sống ở đó, và cả những người vẫn đang ở đó.''

``Vậy thì\dots\ chờ em một chút ạ.''

Tôi bước nhanh đến ô cửa kính lớn ở cuối hành lang, nơi có thể nhìn rõ ra ngoài. Bầu trời bên kia đã hoàn toàn tối sầm, màn đêm bao trùm lấy toàn bộ khu vực văn phòng, chỉ còn những chiếc đèn pha lớn trong bãi đỗ xe vẫn sáng rực, chiếu sáng mọi ngóc ngách.

Tôi dừng lại ở đó, giữ điện thoại sát tai và hạ giọng xuống thấp nhất có thể: ``Chị nói đi ạ, bây giờ chỉ có mình em ở đây thôi. Những văn phòng xung quanh đều đã đóng cửa hết rồi.''

``Kế hoạch này có tên là chiến dịch \ruby{Réunion}{レユニオン},'' giọng của A-san vang lên đều đặn, như thể chị đã chuẩn bị sẵn câu nói này từ lâu.

Tôi im lặng trong vài giây, não bộ phải mất một lúc để xử lý cái tên vừa nghe thấy. Réunion - một từ tiếng Pháp có nghĩa là cuộc đoàn tụ. Tôi lặp lại trong đầu, không dám tin đó là ý nghĩa thực sự của kế hoạch.

``Réunion?''

``Đúng vậy. Mục tiêu của chiến dịch rất đơn giản: tạo ra một cầu nối để một số thành viên hiện tại của công ty có thể gặp lại những người đã tốt nghiệp, nhưng chỉ nếu cả hai phía đều thực sự muốn điều đó. Không có sự ép buộc từ bất kỳ ai.''

Tôi chống một tay lên bậu cửa, lòng bàn tay cảm nhận được sự trơn láng của vật liệu. Câu nói của chị làm tôi phải suy nghĩ rất nhiều. Đây không phải là điều tôi từng nghĩ công ty sẽ đề xuất.

``Vậy nó chỉ là những buổi gặp riêng giữa hai bên? Không phải là một sự kiện công khai ạ?''

``Có thể là một buổi thăm hỏi ngắn, một bữa ăn trưa cùng nhau, hoặc chỉ là một cuộc trò chuyện qua mạng nếu cả hai đều cảm thấy thoải mái. Hình thức cuối cùng sẽ hoàn toàn tùy thuộc vào từng trường hợp, mỗi cặp người sẽ tự quyết định cách họ muốn gặp mặt. Không có một khuôn mẫu nào áp đặt cho tất cả.''

Tôi nuốt nước bọt, một nghi vấn lớn dần hình thành trong đầu, và tôi phải hỏi nó ra ngay: ``Vậy ý chị, nó không phải là một buổi công bố để mời các cựu thành viên quay lại công ty ạ? Không phải là sự kiện tái ký hợp đồng hay gì đó tương tự sao?''

``Hoàn toàn không. Đây không phải là một chiến dịch tuyển dụng mới, không phải lời hứa hay đề nghị bất kỳ ai tái ký hợp đồng với công ty, và cũng không phải là một nội dung để quảng bá trên mạng xã hội hay bất kỳ nền tảng công khai nào cả,'' giọng nói của A-san trở nên dứt khoát hơn bao giờ hết ở câu cuối, như thể chị muốn nhấn mạnh từng chữ để tôi hoàn toàn hiểu rõ quan điểm của công ty.

``Nếu có bất kỳ ai, dù là thành viên hiện tại hay cựu talent, nói rằng họ không muốn gặp mặt, thì chiến dịch sẽ dừng lại ngay lập tức với trường hợp đó. Không có sức ép từ phía công ty, không có thông báo công khai nào về việc ai đã được mời hay ai đã đồng ý, và tuyệt đối không có chuyện dùng tình cảm cũ, những kỷ niệm cũ để buộc ai đó phải đồng ý gặp mặt.''

Tôi thở ra một hơi dài, sự căng thẳng trong người dần vơi đi. Những nguyên tắc mà chị đặt ra đã giải quyết được hầu hết những lo lắng đầu tiên của tôi. Tôi nói nhỏ, tổng kết lại những gì chị vừa trình bày: ``Vậy điều duy nhất mà công ty có thể làm là mở ra một cánh cửa, chứ không có quyền kéo bất kỳ ai phải bước qua cánh cửa đó. Tất cả đều là quyết định của bản thân họ, đúng không ạ?''

``Chính xác. Đó chính xác là những gì chúng ta muốn thực hiện.''

Tôi nhìn xuống con đường chính bên ngoài bãi đỗ xe. Những chiếc xe nối đuôi nhau qua ngã tư gần đó, ánh sáng đỏ từ đèn hậu của chúng tạo thành một dòng dài bất tận, di chuyển chậm chạp trong giờ cao điểm tan sở. Cảnh tượng quen thuộc này làm tâm trạng tôi bình tĩnh lại hơn.

``Vậy ai là người đầu tiên đề xuất ý tưởng này? Ai đã khởi xướng chiến dịch này ạ?'' tôi hỏi.

``YAGOO-san đã từ lâu đặt câu hỏi về khả năng tổ chức những buổi gặp mặt riêng tư như vậy, từ mấy tháng trước rồi. Gần đây, một số talent hiện tại cũng đã chủ động nói với ban quản lý rằng họ muốn có cơ hội nói chuyện lại với những người bạn cũ đã rời công ty, nhưng họ không biết phải bắt đầu từ đâu, làm sao để liên lạc mà không làm phiền đến cuộc sống hiện tại của đối phương.''

Tôi khẽ gật đầu, dù chị không thể thấy. Điều đó hoàn toàn dễ hiểu. Rất nhiều người vẫn giữ liên lạc, nhưng cũng có không ít người mất dấu sau khi hợp đồng kết thúc, vì sợ làm phiền cuộc sống mới của nhau.

``Và\dots\ chị nghĩ rằng ký túc xá là một nơi phù hợp để tổ chức những buổi gặp mặt đầu tiên?''

``Chị không chỉ nghĩ ký túc xá phù hợp. Chị nghĩ em là người phù hợp nhất để làm đầu mối liên kết cho những buổi gặp đó, nếu em đồng ý. Em hiểu rõ tất cả những người đang sống trong nhà, em cũng giữ được mối quan hệ tốt với hầu hết các cựu thành viên. Chị không thể nghĩ ra ai khác phù hợp hơn em.''

Tôi vẫn im lặng, không đưa ra câu trả lời ngay lập tức. Một cơ hội để mọi người có thể gặp lại nhau, những người bạn cũ có thể nối lại mối quan hệ đã đứt đoạn, đó là điều tôi đã mong muốn từ lâu. Nhưng nó cũng đi kèm với quá nhiều trách nhiệm, quá nhiều điều phải cân nhắc để không làm tổn thương bất kỳ ai. Tôi cần thời gian để suy nghĩ kỹ lưỡng.

Từ ngày ký túc xá mở cửa cho các cựu thành viên, căn nhà Hikawa của chúng tôi không chỉ là một nơi ở đơn thuần. mà nó trở thành mái nhà che chở, là vùng an toàn cho những người vừa bước ra khỏi vòng xoáy công việc, cần một khoảng lặng thực sự để làm lại cuộc đời mình sau khi rời khỏi guồng quay của công ty. Nơi này được dựng lên để bảo vệ sự riêng tư, để mọi người được là chính mình mà không phải đối mặt với ánh mắt công chúng hay những áp lực đã từng theo họ suốt thời gian hoạt động.

Nhưng một nơi trú ẩn được xây dựng để che chở không có nghĩa là nó tự động trở thành địa điểm phù hợp cho bất kỳ cuộc gặp gỡ nào, dù mục đích có tốt đẹp đến đâu. Mỗi người trong căn nhà này đều đã xây dựng lại những ranh giới của mình, đã tìm thấy sự ổn định mới sau bao sóng gió, và không ai có quyền phá vỡ sự cân bằng đó một cách tùy tiện.

Những người đang sống ở đây không chỉ là những người ``có mặt'' đơn thuần, họ có quyền được thông báo trước khi có bất kỳ ai bước chân vào ngôi nhà chung này. Họ có quyền biết người đến là ai, lý do cuộc gặp là gì, và liệu nó có ảnh hưởng đến cuộc sống bình yên mà họ đang có hay không.

Quan trọng hơn cả, họ đều có quyền nói lời từ chối. Quyền nói không mà không cần giải thích, không phải chịu bất kỳ sự áp đặt hay trách móc nào. Ngôi nhà này là của họ, và sự thoải mái của họ phải luôn được đặt lên hàng đầu.

``Em muốn biết chị định bắt đầu với ai ạ,'' cuối cùng tôi cũng cất giọng, vẫn giữ điện thoại sát tai, giọng nói nhỏ đến mức chỉ có bản thân mình nghe rõ. Không thể nào đồng ý với một kế hoạch lớn như vậy nếu không biết người đầu tiên bị kéo vào vòng xoáy này là ai, đặc biệt là khi người đó có thể đang sống trong chính căn nhà của tôi.

Đầu dây bên kia vang lên tiếng giấy được lật nhẹ, một tiếng động quen thuộc mà tôi vẫn thường nghe mỗi khi A-san tìm tài liệu trong các cuộc họp. Chị dường như đã chuẩn bị sẵn danh sách, đã lường trước câu hỏi của tôi.

``Nhóm \textit{holoFive},'' giọng chị ấy vang lên, chậm rãi và chắc chắn, như thể đang thăm dò phản ứng của tôi.

Cảm giác như có một lực vô hình vừa bóp nhẹ vào ngực, tim tôi khẽ chùng xuống một nhịp. Mọi lo lắng, mọi dự cảm không tốt bỗng chốc trở thành hiện thực. Tôi đã nghĩ đến nhiều khả năng, nhưng cái tên này vẫn đến bất ngờ, đẩy mọi cảm xúc trong tôi lên đến đỉnh điểm.

``NePoLaBo muốn gặp Aloe-chan ạ?'' câu hỏi tuột ra khỏi miệng tôi trước khi kịp suy nghĩ, mang theo cả sự hoài nghi và không dám tin. Ngay cả khi nghe chị nói thế hệ thứ Năm, tôi vẫn không dám chắc rằng họ muốn gặp chính người mà chúng ta đều phải gọi là Delutaya bây giờ - là Mano Aloe, người chị em cùng nhóm đã rời công ty từ lâu.

``Một vài người trong nhóm đã chủ động hỏi thăm về cô ấy với ban quản lý,'' chị A giải thích, giọng nói nhẹ nhàng như muốn xoa dịu sự căng thẳng của tôi. ``Họ không có yêu cầu phải gặp mặt ngay lập tức đâu. Họ chỉ muốn biết liệu có khả năng nào, có cơ hội nào để có thể nói chuyện lại với Mano-san trong tương lai hay không, nếu cô ấy cũng đồng ý.''

Tôi đưa mắt nhìn vào phản chiếu của chính mình trên cửa kính lớn ở cuối hành lang, dòng chữ mờ ảo của logo công ty in trên mặt kính hiện lên mờ nhạt sau lưng tôi. Trong ánh sáng yếu ớt của hành lang vắng lặng, tôi chỉ thấy một gương mặt đầy lo lắng, đang phải gánh một trách nhiệm nặng nề vừa được đặt lên vai.

``Bây giờ cô ấy không còn là Mano Aloe nữa, chị ạ,'' tôi nói, giọng tôi khẽ run, như đang nhắc nhở chính chị và cả bản thân mình. ``Cô ấy đang sống với cái tên mới, Delutaya. Cuộc đời của cô ấy bây giờ hoàn toàn khác rồi.''

A-san đáp lại ngay lập tức, không có chút bối rối nào: ``Chị hiểu rõ điều đó mà.''

``Cô ấy đang có công việc mới, lịch sinh hoạt ổn định, và có một cuộc sống của riêng mình,'' tôi tiếp tục, liệt kê từng điểm một như muốn nhấn mạnh rằng mọi thứ đã thay đổi, không thể quay lại như xưa. ``Chúng ta không thể cứ mặc định rằng vì từng là bạn bè cùng nhóm, từng trải qua bao nhiêu kỷ niệm, thì mọi chuyện sẽ tự nhiên dễ dàng được đâu.''

``Chị hoàn toàn đồng ý với em,'' chị đáp lại một cách thẳng thắn. ``Chị cũng hiểu rằng mọi ranh giới của cô ấy cần được tôn trọng tuyệt đối.''

``Và có một điều em phải nói rõ,'' tôi nhấn mạnh, giọng nói trở nên cứng rắn hơn, không để có bất kỳ sự hiểu lầm nào. ``Em không bao giờ muốn có ai xem việc này là một kế hoạch để đưa cô ấy trở lại trước công chúng, là bước đầu để cô ấy quay lại với ngành giải trí, hay bất kỳ điều gì tương tự. Đây không phải là một chiến dịch tuyển dụng, và nó không bao giờ được phép trở thành như vậy.''

``Chị đảm bảo sẽ không có chuyện đó xảy ra,'' chị A cam kết, giọng nói đầy chắc chắn. ``Chị sẽ nói rõ điều này với tất cả mọi người liên quan, không để có ai hiểu sai ý nghĩa của cuộc gặp.''

Tôi nghe tiếng chị gõ nhẹ các ngón tay lên mặt bàn ở đầu dây bên kia, một thói quen nhỏ khi chị đang suy nghĩ, đang cân nhắc lời nói của mình.

``Kuon-san, chị gọi em vì chị tin rằng em là người duy nhất hiểu điều đó hơn bất kỳ ai khác,'' chị nói một cách chân thành. ``Em thừa hiểu đây không phải là một cuộc hẹn quay video, không phải một lịch stream có thể sắp xếp một, hai ngày là xong. Đây là chuyện liên quan đến cảm xúc, đến cuộc sống của những người thật, những người đã có quá khứ không dễ dàng, và em là người duy nhất có thể xử lý nó một cách thận trọng.''

Tôi vẫn im lặng, không thể đưa ra bất kỳ câu trả lời nào ngay lúc đó. Trong đầu tôi, vô số hình ảnh ùa về, hiện lên rõ ràng như đang xảy ra trước mắt. Tôi nhớ đến phòng thu nhỏ ở tầng hai của căn nhà, nơi Delutaya đang dành hàng giờ mỗi ngày để sống với đam mê âm nhạc của mình. Tôi nhớ hình ảnh cô ấy ngồi hàng giờ bên chiếc máy tính, chỉnh sửa từng nốt nhạc, từng đoạn hòa âm, để cho ra đời những sản phẩm mà cô ấy tự hào.

Tôi còn nhớ những buổi chiều, cô ấy mang chiếc máy tính xuống phòng khách, bật bản thu mới cho mọi người cùng nghe, rụt rè hỏi ý kiến của mọi người, mắt sáng rực khi nhận được những lời khen. Cái nhìn trong trẻo đó, sự hạnh phúc đơn giản khi được làm điều mình yêu thích, tôi không bao giờ muốn làm tổn thương nó.

Tôi thực sự không biết Delutaya sẽ phản ứng như thế nào, sẽ nghĩ gì khi một ngày nào đó, cái tên NePoLaBo được nhắc lại trước mặt cô ấy. Liệu cô ấy sẽ mừng rỡ, sẽ muốn gặp lại những người chị em cũ? Hay cô ấy sẽ cảm thấy áp lực, sẽ muốn giữ khoảng cách với quá khứ mà cô ấy đã cố gắng để lại phía sau? Không ai có thể đoán trước được.

Và quan trọng nhất, tôi không bao giờ có quyền thay cô ấy đưa ra quyết định đó. Cuộc sống của cô ấy là của cô ấy, sự thoải mái của cô ấy phải là ưu tiên hàng đầu, và chỉ có cô ấy mới có quyền nói rằng mình muốn gì, không muốn gì.

``Em sẽ hỏi Delutaya trước,'' tôi nói.

``Còn nếu như Delutaya hoàn toàn từ chối, không muốn có bất kỳ cuộc gặp nào với họ thì sao?''

``Thì kế hoạch sẽ dừng lại ngay lập tức. Không cần giải thích gì thêm ạ.''

``Còn nếu cô ấy không từ chối hẳn, mà chỉ nói rằng mình cần thêm thời gian để suy nghĩ?''

``Dù cô ấy cần bao lâu, em cũng sẽ đợi. Không có bất kỳ thời hạn nào ép buộc cô ấy phải đưa ra quyết định.''

``Nếu như cô ấy đồng ý gặp mặt, nhưng muốn thay đổi hoàn toàn hình thức mà chị đề xuất ban đầu, ví dụ như gặp ở một nơi khác, hoặc chỉ gặp riêng một mình ai đó thay vì cả nhóm?''

``Thì chúng ta sẽ ngồi lại với cô ấy, thảo luận lại mọi thứ từ đầu. Mọi quyết định cuối cùng đều phải xuất phát từ mong muốn thực sự của cô ấy.''

Đầu dây điện thoại bỗng chìm vào im lặng. Khoảng lặng kéo dài vài giây, chỉ còn lại tiếng hơi thở nhẹ nhàng của chị A xuyên qua loa, như thể chị đang cân nhắc từng chữ trong câu trả lời của mình.

``Đó chính xác là những câu trả lời mà chị đã mong chờ sẽ nhận được từ em đấy, Kuon-san,'' giọng chị A cuối cùng cũng vang lên, mang theo một chút nhẹ nhõm mà tôi có thể cảm nhận rõ ràng.

Lông mày tôi tự nhiên nhíu lại, một chút bối rối len lỏi trong đầu.

``Chị đã nghĩ rằng em sẽ vồn vã đồng ý với kế hoạch này ngay lập tức mà không đặt ra bất kỳ điều kiện nào sao?''

``Không đâu,'' chị lắc đầu, dù tôi không thể thấy cử chỉ đó qua điện thoại. ``Chị đã biết trước rằng em sẽ đặt ra những câu hỏi này. Em luôn là người suy nghĩ thấu đáo cho mọi người hơn bất kỳ ai khác.''

Tôi không khỏi bật cười nhẹ, sự căng thẳng dồn nén trong người từ nãy giờ cuối cùng cũng xua tan đi một phần. Tôi dựa lưng vào khung cửa kính lạnh lẽo, cảm nhận được hơi ẩm của đêm mùa thu se lạnh thấm qua lớp áo.

``Vậy thì, cuộc gọi này có phải là cách chính thức để thông báo rằng chiến dịch Réunion chính thức bắt đầu chạy rồi không?''

``Không hẳn là bắt đầu đâu,'' chị A sửa lời nhẹ nhàng. ``Đúng hơn là chị muốn báo trước cho em biết rằng kế hoạch này đã được ban lãnh đạo duyệt về mặt nguyên tắc. Tất cả những gì chúng ta có bây giờ chỉ là một khung sườn trống. Mỗi cuộc gặp riêng sau này đều phải được sự đồng thuận hoàn toàn từ cả hai phía, không có ngoại lệ.''

``Vậy ai sẽ là người trực tiếp liên lạc với các bạn của NePoLaBo? Chị hay là có ai khác phụ trách mảng đó?''

``Chị sẽ tự mình làm việc này. Em không cần phải lo lắng về khía cạnh đó. Nhưng trước khi chúng ta làm bất cứ điều gì với họ, điều tiên quyết là phải biết được Delutaya-san có muốn nghe về lời đề nghị này hay không. Nếu cô ấy không muốn, chúng ta sẽ không bao giờ đề cập cái tên này trước mặt cô ấy.''

``Tuyệt đối đừng gửi bất kỳ lời mời chính thức nào, đừng có bất kỳ văn bản hay thông báo nào được gửi đến cô ấy trước khi em có cơ hội nói chuyện riêng với cô ấy,'' tôi nhấn mạnh từng chữ, giọng tôi trở nên cứng rắn hơn một chút. Tôi cần phải chắc chắn rằng chị A hiểu rõ tầm quan trọng của điều này.

``Chị sẽ không bao giờ làm vậy. Em hoàn toàn có thể yên tâm.''

``Và với cả, xin chị đừng bao giờ nhắc bất cứ điều gì liên quan đến kế hoạch này trước mặt những người còn lại đang sống trong ký túc xá. Chưa có gì chắc chắn cả, và em không muốn tạo ra bất kỳ kỳ vọng không cần thiết nào cho ai hết.''

``Chị hiểu rõ rồi. Chị sẽ giữ kín chuyện này hoàn toàn. Em hãy tự mình lựa chọn thời điểm thích hợp nhất để nói chuyện với Delutaya, không cần phải vội vàng.''

Tôi vô tình liếc nhìn xuống chiếc đồng hồ đeo tay, kim đồng hồ đang nhích dần từng giây. Cơn sốc ban đầu khi nghe chị A đề cập kế hoạch, những lo lắng, những câu hỏi xoay chuyển trong đầu—tất cả những thứ đó khiến tôi cảm giác như đã trải qua hàng giờ dài. Nhưng trên thực tế, cuộc gọi này mới chỉ kéo dài chưa đầy mười phút.

Thế nhưng, trong lòng tôi lại cảm giác như thể một cánh cửa cũ kỹ, bị đóng chặt từ rất lâu ở cuối một hành lang tối tăm và dài vô tận, vừa mới được hé mở. Một chút ánh sáng yếu ớt từ phía bên kia lọt vào, đủ để khiến con tim tôi tràn đầy hy vọng, nhưng cũng đủ để tôi nhận ra rằng con đường phía trước vẫn còn rất nhiều điều chưa thể biết trước.

\ct

``Chị còn một điều cuối cùng muốn nhắc em nhớ,'' giọng chị A bỗng kéo tôi trở lại từ dòng suy nghĩ miên man.

``Vâng ạ, chị nói đi.''

``Khi em nói chuyện với Delutaya, đừng bao giờ gọi đây là 'chiến dịch'. Đừng dùng cái tên mà chúng ta vẫn dùng trong các cuộc họp nội bộ.''

Tôi hơi ngẩn người ra, một chút ngạc nhiên trước yêu cầu bất ngờ đó. Lúc đầu tôi không hiểu tại sao lại có sự khác biệt nhỏ như vậy, nhưng sau đó tôi dần nhận ra ý đồ đằng sau nó.

``Tại sao lại vậy ạ? Có vấn đề gì nếu em gọi nó là chiến dịch không?''

``Bởi vì cái tên 'chiến dịch' chỉ dành cho những cuộc thảo luận, những sự phối hợp nội bộ giữa chúng ta trong công ty. Đó là một thuật ngữ mang hơi hướng công việc, có tính tổ chức. Còn đối với Delutaya, đây không bao giờ được phép là một dự án của công ty. Với cô ấy, đây chỉ nên đơn thuần là một lời hỏi thăm chân thành từ những người bạn cũ, và một lựa chọn hoàn toàn tự do mà cô ấy có thể chấp nhận hoặc từ chối mà không cần bất kỳ lý do nào. Chỉ có vậy thôi.''

Tôi nhanh chóng gật đầu, dù chị A không thể thấy được hành động đó qua màn hình điện thoại. Tôi hoàn toàn đồng ý với suy nghĩ của chị: ``Em sẽ nhớ kỹ điều đó. Không bao giờ nhắc đến hai từ 'chiến dịch' khi nói chuyện với cô ấy.''

``Cảm ơn em đã hiểu chuyện, Kuon-san. Chị thực sự biết ơn vì em đã đồng ý giúp đỡ chúng ta trong việc này.''

``A-san này, em còn có một vài điều muốn nhắc chị lần nữa để đảm bảo rằng chúng ta đều thống nhất.''

``Ừ, em nói đi. Bất cứ điều gì em muốn nhắc, chị đều lắng nghe.''

``Đừng để bất kỳ ai sử dụng hình ảnh, tin nhắn hay thậm chí là sự đồng ý của một người để làm áp lực buộc người còn lại phải đồng ý. Không ai có quyền dùng tình cảm cũ, hay bất cứ điều gì tương tự để ép buộc người khác phải làm những điều họ không muốn.''

``Chị sẽ viết điều đó vào những nguyên tắc cốt lõi của chiến dịch. Tất cả mọi người tham gia đều phải đọc và ký xác nhận rằng họ hiểu rõ điều này.''

``Và một điều nữa, tuyệt đối không được phép biến bất kỳ cuộc gặp riêng nào thành nội dung công khai trên mạng xã hội hay bất kỳ nền tảng nào, nếu chúng ta chưa hỏi ý kiến và nhận được sự đồng ý từ cả hai bên. Tất cả đều phải là riêng tư, hoàn toàn kín đáo.''

``Chị cam kết với em. Những điều em đề cập sẽ là những quy tắc bất di bất dịch, không ai được phép vi phạm.''

Tôi tin tưởng vào lời cam kết của chị A. Tôi biết chị luôn giữ lời nói của mình, và những nguyên tắc này sẽ được chị bảo vệ nghiêm ngặt.

Nhưng niềm tin vào chị không có nghĩa là tôi có thể bỏ qua việc phải nói rõ những ranh giới này ra. Niềm tin không bao giờ thay thế được sự minh bạch, những quy tắc được đặt ra ngay từ đầu để tránh mọi hiểu lầm có thể xảy ra trong tương lai.

Tôi cần phải nói hết tất cả những điều này ra, để mọi thứ đều rõ ràng, trước khi bất kỳ ai có cơ hội bước chân vào căn nhà Hikawa của chúng ta. Căn nhà mà tôi đã dành rất nhiều tâm sức để biến nó thành một nơi an toàn cho tất cả những người mà tôi quan tâm.

``Vậy thì, em sẽ đồng ý làm đầu mối liên kết cho buổi gặp đầu tiên này,'' tôi cuối cùng cũng đưa ra quyết định, giọng tôi chắc chắn và vững vàng.

``Nhưng chỉ khi Delutaya đồng ý với mọi thứ thôi.''

``Chỉ khi Delutaya thực sự muốn cuộc gặp này diễn ra.''

Tiếng tút tút ngắt quãng vang lên từ đầu dây khi cuộc gọi chính thức kết thúc, để lại khoảng lặng chỉ có tiếng hơi thở nhẹ nhàng của chính mình vang vọng trong hành lang vắng lặng. Tôi vẫn giữ nguyên tư thế đó, vẫn tựa vai vào khung cửa kính lạnh lẽo, ánh mắt vẫn hướng ra bãi đỗ xe trống trải phía ngoài mà không thực sự nhìn vào bất cứ điều gì cụ thể. Những dòng cảm xúc còn đang xoáy chặt trong lòng, chưa thể nào bình tĩnh lại ngay lập tức.

Thời gian trôi qua từng giây dài mà tôi chẳng hề hay biết, cho đến khi một tiếng rung nhẹ vang lên từ chiếc điện thoại vẫn còn cầm trong tay, một xung động nhỏ xuyên qua lòng bàn tay, kéo tôi trở lại với hiện tại.

Tôi nâng màn hình lên, dòng chữ hiện rõ ràng dưới tên người gửi là A-san:

\txtbx{``Cảm ơn Kuon-san đã đồng ý hợp tác. Chưa có gì được sắp xếp hay lên lịch cả, mọi thứ sẽ dừng lại hoàn toàn cho đến khi Delutaya-san tự mình đưa ra quyết định.''}

Ngón tay tôi lướt nhẹ trên màn hình cảm ứng, gõ ra một câu trả lời ngắn gọn:

\txtbx{``Đã rõ. Em sẽ tìm thời điểm phù hợp nhất để nói chuyện riêng với cô ấy.''}

Nhấn nút gửi, tôi vừa định đặt điện thoại xuống thì giọng nói quen thuộc của Game vang lên từ chiếc tai nghe vẫn còn đeo trên tai, cất lên nhẹ nhàng trong không khí im lặng: ``Tôi đã ghi nhận toàn bộ nội dung cuộc gọi vừa rồi.''

``Đừng lưu bất kỳ bản ghi âm nào của cuộc gọi này,'' tôi lập tức nhắc nhở, giọng vẫn còn giữ ở mức thấp, như thể sợ rằng có ai đó có thể tình cờ nghe thấy. Điều về sự riêng tư này tôi không bao giờ coi thường.

``Tôi không ghi âm cuộc gọi. Tôi chỉ ghi lại những ghi chú công việc cần thiết, không lưu lại bất kỳ đoạn âm thanh nào.''

``Vậy thì được rồi,'' tôi thở nhẹ ra một hơi, phần nào yên tâm hơn.

``Cậu có vẻ đang rất lo lắng về việc này đấy, Kuon-user,'' Game nhận xét, giọng nói mang theo sự quan sát tinh tế mà nó luôn có.

``Tôi không muốn một cuộc gặp gỡ riêng tư, vốn chỉ nên xuất phát từ mong muốn chân thành của cả hai bên, lại vô tình biến thành một nghĩa vụ mà ai đó phải gánh vác,'' tôi giải thích, những lo lắng trong lòng được nói ra thành lời, cũng nhẹ nhõm đi một phần. ``Tôi cũng không muốn bất kỳ ai cảm thấy bị ép buộc phải tham gia chỉ vì những tình cảm cũ, hay những mối quan hệ đã từng có.''

``Đó chính xác là lý do tại sao cậu lại đặt ra những điều kiện rõ ràng ngay từ đầu, phải không? Để ngăn chặn mọi khả năng đó xảy ra.''

Tôi cất chiếc điện thoại vào sâu trong túi quần, như thể đang cất đi cả gánh nặng nặng nề vừa mới đặt lên vai này. Trở về với bàn làm việc của mình, những công việc thường nhật vẫn chất đống đó, các cuộc họp nhỏ lẻ vẫn diễn ra đúng lịch trình, ngày làm việc tiếp tục trôi đi như bao ngày bình thường khác, chẳng có gì thay đổi trên bề mặt.

Nhưng trong sâu thẳm tâm trí tôi, cụm từ ``Réunion'' cứ lặp đi lặp lại không ngừng, như một bản nhạc không lời cứ vong vọng mãi không thôi. Nó không hiện lên như cái tên của một kế hoạch công ty, một chiến dịch được lập trình kỹ lưỡng với các mục tiêu và thời hạn rõ ràng. Nó chỉ đơn thuần hiện lên như một lời gõ cửa nhẹ nhàng, một tiếng động nhỏ ở đầu một hành lang dài tối tăm, như có ai đó đang đứng đợi bên ngoài, chờ đợi xem liệu cánh cửa cũ kỹ kia có được hé mở hay không.

\ct

Cái đêm sau cuộc gọi với A-san, tôi bước chân vào căn nhà Hikawa sau bao nhiêu công việc còn dang dở ở công ty, thời gian đã trôi qua quá nhanh, khiến tôi về đến nhà muộn hơn mọi ngày thường. Ánh đèn vàng ấm áp ở hành lang đón tôi, nhưng cả tầng năm bỗng chốc trở nên yên tĩnh lạ thường: không còn tiếng cười đùa, không còn tiếng máy chiếu rền rĩ như những buổi tối xem phim cùng nhau.

Tôi bước nhẹ vào phòng khách, thấy chiếc hộp gỗ đựng máy chiếu của Kson đã được đóng kín, đặt gọn gàng trên kệ tivi, chứng tỏ buổi xem phim tối nay đã kết thúc từ lâu.

Căn phòng Naomi và Koneko ở dưới tầng ba, cửa đã đóng chặt, ánh đèn nhỏ trên bàn cạnh giường đã tắt; hai đứa chắc hẳn đã ngủ say từ lúc nào. Căn phòng của Ryuuhou còn lọt ra một chút ánh sáng yếu ớt, nghe thấy tiếng lật sách nhẹ, còn Suzuno đang đứng ở bếp nhỏ, dọn dẹp mấy chiếc cốc thủy tinh còn lại trên bàn, lau sạch vết trà còn sót lại.

Đi ngang qua tầng hai, tôi dừng chân trước cửa phòng thu. Dù được lắp lớp kính cách âm dày đặc, nhưng vẫn có thể nghe thấy một đoạn nhạc nhỏ thoát ra, nhẹ nhàng như tiếng mưa rơi trên mái nhà. Delutaya đang ở trong đó, như thường lệ, dồn toàn bộ tâm trí vào bản nhạc mới cô ấy đang hoàn thiện, chỉnh sửa từng nốt nhạc từng đoạn hòa âm cho đến khi hài lòng.

Tôi đứng yên ngoài cửa một lúc lâu, bàn tay đã đặt trên tay nắm nhưng lại không dám xoay mở. Không muốn làm phiền, không muốn làm gián đoạn khoảnh khắc cô ấy đang hòa mình vào âm nhạc, khoảnh khắc mà cô ấy có thể hoàn toàn là chính mình, không có bất kỳ áp lực nào đè nặng.

Tôi định rời đi, để lại chuyện đó cho ngày mai khi mọi thứ đều đủ bình tĩnh, khi tôi có thể nói chuyện với cô ấy một cách nhẹ nhàng nhất. Nhưng đúng lúc tôi vừa quay lưng, chiếc điện thoại trong túi quần bỗng rung lên, một thông báo tin nhắn hiện lên màn hình.

Tôi mở ra, thấy nhóm chat công việc bật lên một tin nhắn mới từ Nene:

\txtbx{``Kuon-senpai, cuối tuần này bọn em có thể ghé thăm nhà anh được không? Lamy-chan với Botan-chan cũng muốn sang, chúng em muốn ghé chơi thôi ạ!''}

Tôi đọc đi đọc lại dòng chữ đó mấy lần, như thể không dám tin vào mắt mình. Mới chỉ vài tiếng trước, tôi và A-san vừa bàn về kế hoạch Réunion, vừa nói rằng mọi thứ còn chưa có gì chắc chắn, rằng chúng ta phải đợi Delutaya đồng ý trước khi làm bất cứ điều gì. Chưa có một lời mời nào được gửi đến cô ấy, chưa có một buổi gặp nào được lên lịch hay xác nhận, nhưng như thể có một sức mạnh vô hình nào đó đang đẩy mọi thứ tiến về phía trước: cánh cửa đầu tiên mà chúng ta còn đang đắn đo có nên mở hay không, bỗng chốc tự nó đã hé mở một chút.

Tôi gõ trả lời Nene, giọng điệu giữ vững sự bình thường:

\txtbx{``Để anh kiểm tra lịch của mọi người trong nhà đã, rồi sẽ báo lại em sớm.''}

Tôi không hề nhắc gì đến cuộc gọi của A-san mấy tiếng trước, cũng không hề hé lộ rằng đây có thể là bước đầu tiên của chiến dịch mà chúng ta đã chuẩn bị bao lâu. Trước hết, Delutaya phải được nghe chuyện này từ chính miệng tôi. Và quan trọng nhất, cô ấy phải có toàn quyền lựa chọn: muốn hay không, gặp hay không, tất cả đều phụ thuộc vào mong muốn của cô ấy, không có bất kỳ sức ép nào.

%==========================================TRUYỆN PHỤ=========================================
\omk

Sáng hôm sau, khi ánh nắng mùa thu vàng ấm vừa len lỏi qua cửa sổ tầng hai, tôi lại đứng trước cánh cửa phòng thu đó. Tiếng nhạc nhỏ vẫn vọng ra từ bên trong, nhè nhẹ, đều đặn; Delutaya đã sớm vào làm việc, như mỗi ngày.

Tôi đưa tay lên, các ngón tay đã chạm vào bề mặt cửa gỗ, định gõ nhẹ để báo hiệu có người ngoài, nhưng rồi lại dừng lại. Không có lý do gì phải vội vàng cả, không có bất kỳ thời hạn nào ép buộc tôi phải nói chuyện với cô ấy ngay lập tức. Tôi hạ tay xuống, đứng yên ngoài hành lang, chờ đợi.

Tôi sẽ đợi đến khi cô ấy nghỉ ngơi, khi cô ấy bước ra khỏi phòng thu với tâm trạng thoải mái, rồi mới hỏi chuyện.

Tôi không muốn đứng trước mặt cô ấy với tư cách là người quản lý của căn nhà, hay là người đại diện cho công ty đưa ra một kế hoạch đã được ban lãnh đạo duyệt qua. Tôi chỉ muốn đến với cô ấy với tư cách một người bạn - muốn hỏi rằng, ``liệu em có muốn gặp lại những người bạn cũ của mình không?'' với những người đã từng cùng cô nàng trải qua bao nhiêu kỷ niệm, những người em vẫn thường nhắc đến trong những buổi trà chiều thong thả.

Rồi tiếng nhạc trong phòng bỗng dừng lại. Một lát sau, ổ khóa cửa xoay nhẹ, cánh cửa mở ra một khoảng nhỏ, Delutaya ló đầu ra ngoài, dây tai nghe vẫn quàng quanh cổ, mắt còn đang lưu luyến với màn hình máy tính trong phòng. Em ấy ngạc nhiên khi thấy tôi đứng ngoài, khẽ nghiêng đầu: ``Kuon-kun? Anh đứng đây từ lúc nào ạ? Anh cần em giúp gì không?''

Tôi nhìn cô ấy, rồi mỉm cười nhẹ nhàng, cố gắng để giọng nói của mình thật bình thường, không mang theo bất kỳ sự nặng nề nào: ``Không có gì gấp đâu. Khi nào em nghỉ xong hết công việc, anh muốn hỏi em một chuyện, không quan trọng lắm.''

Delutaya lập tức gật đầu, nụ cười thân thương nở trên môi: ``Vâng ạ, em chỉ còn chỉnh sửa nốt một đoạn nhỏ thôi. Anh đợi em một chút nhé, em ra ngay.''

Tôi lùi sang bên, nhường lối cho cô ấy nếu cần ra ngoài, hay để cô ấy đóng cửa lại tiếp tục công việc. Trong lòng tôi chỉ có một quyết định: lần này, tôi sẽ không bao giờ thay cô ấy đưa ra bất kỳ lựa chọn nào. Tất cả mọi thứ đều phụ thuộc vào cô ấy, và nếu câu trả lời cuối cùng là không - nếu cô ấy không muốn gặp lại những người bạn cũ, không muốn mở lại cánh cửa quá khứ đó - thì cánh cửa đó vẫn sẽ ở nguyên đó, vẫn sẽ được đóng chặt, không có ai có quyền ép buộc cô ấy phải mở nó ra.

Mặt trời dần lặn xuống phía sau những nhà cao tầng xa xăm, nhuốm màu cam ấm áp lên bầu trời chiều, khi chiếc điện thoại trong túi quần tôi lại rung lên. Là chị A, gọi lại như đã hứa từ sáng sớm. Tôi bước ra khỏi phòng làm việc nhỏ của mình ở tầng năm, tránh xa những tiếng cười đùa thoảng qua từ phòng khách, nơi Suzuno và Ryuuhou đang ngồi xem chương trình truyền hình, để có một không gian riêng tư để nghe điện thoại.

``Em đã nói chuyện với Delutaya về việc đó chưa, Kuon-san?'' giọng nói của chị A vang lên từ đầu dây, mang theo sự quan tâm nhẹ nhàng, không có chút vội vã nào.

Tôi dựa lưng vào bức tường gỗ của hành lang, mắt vẫn hướng về cánh cửa phòng thu ở cuối hành lang, nơi âm nhạc nhỏ vẫn tiếp tục lan tỏa ra: ``Em mới chỉ nói với cô ấy rằng có một chuyện muốn hỏi, chưa trình bày cụ thể nội dung gì hết. Em muốn đợi cô ấy sẵn sàng, muốn để em ấy tự cảm thấy thoải mái trước khi nghe những chuyện này.''

``Được rồi. Cứ để cô ấy chọn thời điểm mình muốn nghe, không có gì phải vội vàng cả,'' chị A đáp lại, giọng đầy đồng tình. Tôi có thể nghe thấy tiếng lật giấy nhẹ nhàng từ đầu dây bên kia, hình như chị đang xem qua vài tài liệu trên bàn làm việc.

Sau một lúc im lặng nhỏ, chị tiếp tục, giọng nói có chút thay đổi, như thể đang nhắc đến một chuyện mà chị cũng vừa mới nhận ra: ``Còn có một chuyện nữa, Momosuzu-san vừa nhắn tin cho chị, hỏi rằng cuối tuần này nhóm em ấy có thể ghé đến nhà em chơi không.''

Tôi lập tức rút chiếc điện thoại từ trong túi ra, mở màn hình lên. Cái tin nhắn của Nene mà tôi đã đọc từ tối qua vẫn đang nằm nguyên trên màn hình. Tôi nhìn chằm chằm vào dòng chữ đó một lúc rồi mới trả lời chị: ``Em cũng vừa nhận được tin nhắn đó của Nene-chan từ tối qua rồi ạ. Em cũng đang định bàn với chị về việc này.''

``Chị chưa trả lời gì với các em ấy đâu. Chị không muốn xác nhận hay hứa hẹn bất cứ điều gì trước khi em và em ấy đều sẵn sàng. Chị chỉ muốn thông báo trước cho em biết, để em không bị bất ngờ nếu cô ấy có nhắc lại chuyện đó với em,'' chị A nói thật thà, rõ ràng chị cũng đang rất thận trọng với từng bước đi của mình.

``Cảm ơn chị đã luôn thông báo trước cho em như vậy, em rất biết ơn,'' tôi nói, lòng cảm thấy nhẹ nhõm đi một phần vì chị A cũng đang giữ cùng một nguyên tắc với mình, không để bất cứ điều gì xảy ra một cách vội vàng.

Rồi giọng chị A chậm rãi vang lên, như thể nhắc lại một nguyên tắc cốt lõi mà chúng ta đã thống nhất từ đầu: ``Kuon-san, chị muốn nhắc lại một lần nữa, chiến dịch Réunion của chúng ta chỉ có thể chính thức bắt đầu khi tất cả mọi người đều tự bước tới. Không ai bị ép buộc, không ai bị kéo theo, mọi thứ đều phải đến từ mong muốn thực sự của từng người.''

``Em hiểu rõ điều đó, chị ạ. Em sẽ không bao giờ để bất cứ ai phải cảm thấy bị ép buộc phải làm gì mình không muốn,'' tôi đáp lại chắc chắn, những lời đó cũng chính là những gì tôi luôn nhắc nhở bản thân mình.

Một lát sau, cuộc gọi kết thúc với những lời chúc bình an từ hai phía. Tôi cất chiếc điện thoại trở lại vào trong túi quần, đứng yên một mình trong hành lang yên tĩnh, mắt vẫn hướng về phía cánh cửa phòng thu cách đó vài bước chân.

Từ bên trong căn phòng nhỏ đó, tôi vẫn có thể nghe thấy giọng của Delutaya đang thử lại một đoạn điệp khúc của bản nhạc mới, em ấy lặp đi lặp lại từng câu từng chữ để tìm ra giọng điệu ưng ý nhất, từng nốt nhạc lan tỏa ra nhẹ nhàng, tràn đầy sự tập trung và đam mê. Tôi đứng yên ngoài cửa, bàn tay lại một lần nữa chạm vào tay nắm cửa gỗ ấm áp, nhưng rồi lại rút ra.

Lần này tôi không gõ cửa. Tôi không muốn làm phiền khoảnh khắc em ấy đang hòa mình hoàn toàn vào âm nhạc, không muốn phá vỡ sự tập trung đó. Tôi sẽ đứng đây đợi, đợi cho đến khi em ấy nghỉ xong tất cả công việc, đợi cho đến khi em ấy bước ra khỏi phòng với nụ cười thoải mái, rồi mới sẽ nói chuyện với em ấy về tất cả những chuyện này. Mọi thứ đều có thể chờ đợi, không có gì là quan trọng hơn sự thoải mái và tự do lựa chọn của cô nàng Tam Giác này.

\end{document}